\documentclass[11pt,letterpaper]{article}
\usepackage{cornell-notes}
\usepackage{mathtools}

\title{Chapter 1: Exact Matching: Fundamental Preprocessing and First Algorithms}
\author{}
\date{}

\setDocTitle{Chapter 1: Exact Matching: Fundamental Preprocessing and First Algorithms}
\setDocSubtitle{String Algorithms}
\setDocVersion{v1.0}
\setDocStatus{Study Companion}
\setDocOwner{String Algorithms Cornell Notes}
\setDocAuthor{}
\setDocDate{\today}
\setDocSummary{Cornell-format study and review notes for Chapter 1: Exact Matching: Fundamental Preprocessing and First Algorithms.}

\setCornellCollection{String Algorithms}
\setCornellUnitType{Chapter}
\setCornellUnitNumber{1}
\setCornellUnitTitle{Exact Matching: Fundamental Preprocessing and First Algorithms}
\setCornellCompanionLabel{Study and review companion}

\hypersetup{pdftitle={Chapter 1: Exact Matching: Fundamental Preprocessing and First Algorithms},pdfsubject={Cornell notes},pdfauthor={}}

\begin{document}
\maketitle

\begin{CornellOverviewBox}{Chapter overview}
\textbf{Purpose.} Develop the preprocessing viewpoint for exact pattern matching and derive a linear-time search directly from prefix-match information.

\textbf{Learning objectives}
\begin{itemize}
  \item Distinguish naive alignment testing from preprocessing-based search.
  \item Define $Z_i$, a $Z$-box, and the maintained interval $[\ell,r]$.
  \item Explain why the $Z$ algorithm performs only linear total comparison work.
  \item Use $P\$T$ to find all exact occurrences of a pattern.
\end{itemize}

\textbf{Prerequisites.} Strings, prefixes, substrings, asymptotic notation, and array indexing.
\end{CornellOverviewBox}

\begin{CornellNotesTable}
\CornellNoteRow{\S1.1: What does the naive method do?}{Align a pattern $P$ of length $m$ at each feasible position of a text $T$ of length $n$ and compare corresponding characters until a mismatch or a full match. In the worst case it performs $\Theta(m(n-m+1))$ comparisons. Its weakness is not correctness but failure to reuse information learned at earlier alignments.}
\CornellNoteRow{How can comparisons be skipped safely?}{A mismatch after a long partial match reveals structure about the aligned text segment. If the pattern has been preprocessed, the search can shift to the next alignment consistent with that structure and can sometimes avoid rechecking characters already known to agree.}
\CornellNoteRow{\S1.2: What is preprocessing?}{Separate the work into a preprocessing stage and a search stage. Pattern preprocessing records internal repetitions before the text is known; text preprocessing, used by later indexing structures, supports many future pattern queries. The useful split depends on which string is reused.}
\CornellNoteRow{\S1.3: What is $Z_i$?}{For a string $S$, $Z_i$ is the length of the longest substring beginning at position $i>1$ that equals a prefix of $S$. If $Z_i>0$, the interval $[i,i+Z_i-1]$ is a $Z$-box. These values summarize every prefix occurrence in $S$.}
\CornellNoteRow{What are $\ell$ and $r$?}{While scanning positions left to right, $[\ell,r]$ denotes a previously discovered $Z$-box whose right endpoint is farthest right. The invariant is that $S[\ell..r]=S[1..r-\ell+1]$. A position inside this box can borrow a previously computed $Z$ value from its corresponding prefix position.}
\CornellNoteRow{\S1.4: How is $Z_k$ computed?}{If $k>r$, compare explicitly from $S[k]$ and the prefix. If $k\le r$, let $k'=k-\ell+1$. When $Z_{k'}<r-k+1$, set $Z_k=Z_{k'}$ without new comparisons. Otherwise, the known match reaches $r$; compare only beyond $r$ and update $[\ell,r]$ if the match extends.}
\CornellNoteRow{Why is the $Z$ algorithm linear?}{Each iteration contributes at most one terminal mismatch. Every successful comparison made explicitly advances the global right boundary $r$, which never retreats and never exceeds $|S|$. Thus the total number of matches and mismatches is $O(|S|)$, with $O(|S|)$ stored $Z$ values.}
\CornellNoteRow{\S1.5: How does $P\$T$ solve exact matching?}{Choose a delimiter $\$$ absent from both strings and form $S=P\$T$. Compute all $Z$ values. A text position is an occurrence exactly when the corresponding $Z_i$ equals $|P|$. Construction and scanning take $O(|P|+|T|)$ time; storing only the pattern-side values reduces auxiliary space to $O(|P|)$.}
\CornellNoteRow{What assumptions matter?}{The method needs only equality tests, so its linear bound is alphabet independent. The delimiter must be outside the working alphabet, and the implementation must translate indices between $S$ and $T$ carefully. Empty-pattern behavior should be specified by the calling API.}
\CornellNoteRow{\S1.6: What should practice emphasize?}{Derive linear reductions to exact matching, such as detecting a cyclic rotation by searching one equal-length string inside a doubled copy of the other. Practice proving both directions of each reduction and accounting for delimiter and boundary cases.}
\end{CornellNotesTable}

\begin{CornellSummaryBox}{Chapter synthesis}
Prefix-match preprocessing turns repeated work into reusable structure. The $Z$ algorithm maintains the rightmost verified prefix copy, reuses earlier $Z$ values inside that interval, and compares explicitly only when it can push the frontier. Applying the same computation to $P\$T$ yields a simple alphabet-independent linear-time matcher.
\end{CornellSummaryBox}

\begin{CornellSummaryBox}{Key takeaways}
\begin{CornellRecallList}
  \item Naive matching forgets partial-match information.
  \item $Z_i$ measures the prefix match beginning at position $i$.
  \item The right boundary $r$ is the key amortization device.
  \item Positions inside a $Z$-box can often be resolved without comparisons.
  \item $P\$T$ converts exact matching into prefix-occurrence detection.
  \item Preprocessing choice should follow which input will be reused.
\end{CornellRecallList}
\end{CornellSummaryBox}

\begin{CornellOverviewBox}{Algorithm selection: when to use this}
\begin{CornellChecklist}
  \item Use the direct $Z$ matcher for a compact linear-time exact search.
  \item Use pattern preprocessing when one pattern is searched in one or more streams.
  \item Prefer a text index when the same large text will receive many queries.
\end{CornellChecklist}
\end{CornellOverviewBox}

\begin{CornellExamBox}{Self-test questions}
\begin{CornellRecallList}
  \item State the $Z_i$ definition without referring to an example.
  \item Why does $Z_{k'}<r-k+1$ determine $Z_k$ immediately?
  \item Give the amortized argument that bounds explicit successful comparisons.
  \item How is a match position in $P\$T$ translated back to $T$?
  \item Why must the delimiter be absent from both inputs?
\end{CornellRecallList}
\end{CornellExamBox}

{\footnotesize
\textbf{Terms and notation to remember.} $Z_i$, $Z$-box, $[\ell,r]$, prefix, preprocessing stage, search stage, delimiter, alphabet independence.

\textbf{Connections.} The same prefix structure is repackaged by classical failure functions, while suffix trees reverse the preprocessing role and index the text.
}

\end{document}
